\section{Tails and mean length consequences}\label{c4:tails}

We now have both projection losses: the unrooted loss from the
contractible pressure and the length-weighted loss from the one-mark
ratio. The transverse-height estimate converts these cylinder
quantities into planar diameter tails.

\subsection{From projection losses to diameter tails}

Return to simple unoriented plane polygons modulo translations by
$\mathbb Zh+\mathbb Zv$, with critical weight
$w(\gamma)=x_*^{|\gamma|}$. Write
\[
 A(R)=\sum_{\diam\gamma\ge R}w(\gamma),\qquad
 P(R)=\sum_{\diam\gamma\ge R}|\gamma|w(\gamma),
\]
where diameter is Euclidean. We prove
\begin{equation}
 A(R)=R^{-2+o(1)},\qquad P(R)=R^{-2/3+o(1)}.               \label{c4:Tail}
\end{equation}

The length-weighted plane mass lost under projection to period $Nh$
is exactly $E_N$. Indeed, the minus-twist marked identity gives total
mass $c_*$ through the three sampled cylinder edges. Its contractible
part is
\[
 \sum_{\text{surviving plane classes }\gamma}|\gamma|w(\gamma).
\]
This is the same marking and translation count as in
\eqref{c4:Press-6}: each surviving plane class has $N$ axial cylinder
classes, while fixing one edge of each type counts its edges and removes
that factor $N$. Each plane polygon eventually survives. Fatou's lemma
therefore gives a finite full plane length-weighted mass, bounded by
$c_*$. Dominated convergence then identifies that mass as $c_*$, since
$E_N\to0$ by Proposition~\ref{prop:one-mark-ratio}. Subtracting the
survivor sum proves the loss identity. The two losses are consequently
\[
 \sum_{\text{not surviving}}w(\gamma)\asymp N^{-2},\qquad
 \sum_{\text{not surviving}}|\gamma|w(\gamma)
       =E_N=N^{-2/3+o(1)},
\]
by \eqref{c4:Press-8} and Proposition~\ref{prop:one-mark-ratio}.

A failed projection has a collision between vertices differing by a
nonzero multiple of $Nh$, so its diameter is at least $|h|N$. This
gives both lower tail bounds. Conversely, every polygon of diameter at
least $R$ has transverse row span at least $cR$ in one of three
symmetry-related orientations of the horizontal period. Fix a small
$\epsilon>0$ and take $N=\lfloor R^{1-\epsilon}\rfloor$ in each
orientation. In the orientation with large span, the polygon is either
lost or projects to a tall contractible cylinder polygon.
Proposition~\ref{prop:thin-cylinder-span} makes the latter mass smaller
than every fixed inverse power of $R$, also with one length weight.
The lattice rotations preserve the translation-class sums. Adding the
three loss bounds and then taking $\epsilon$ arbitrarily small proves
\eqref{c4:Tail}.

\begin{corollary}[Mean polygon lengths]\label{cor:polygon-means}
The conditional mean length in the unrooted critical polygon ensemble
on diameter $\ge R$ is $R^{4/3+o(1)}$. The same is true on
$R\le\diam\gamma<R^{1+\epsilon}$ for every fixed $\epsilon>0$.
The mean length of a critical cylinder winding polygon sampled by
axial classes is likewise $NE_N/A_N=N^{4/3+o(1)}$.
\end{corollary}

\begin{proof}
Divide the length-weighted mass by the unrooted mass in
\eqref{c4:Tail}. Subtracting the larger-diameter tails gives the
power-separated annulus. For cylinder windings, use $NE_N/A_N$ with
\eqref{c4:Press-6} and Proposition~\ref{prop:one-mark-ratio}.
\end{proof}

\subsection{A bounded modification at a lowest vertex}

\begin{corollary}[A half-plane mean with fixed nearby endpoints]
\label{cor:boundary-neighbor-mean}
Take the upper half-plane above the row cut through the midpoints of
the vertical bonds below the $P_{i,0}$. Consider all SAWs there from
$P_{0,0}$ to $P_{1,0}$, weighted by $x_*^\ell$, where $\ell$ is edge
length. Under this normalized weight, conditional on diameter at least
$R$, the mean length is $R^{4/3+o(1)}$ as $R\to\infty$. The same
conclusion holds on $R\le\diam<R^{1+\epsilon}$ for every fixed
$\epsilon>0$.
\end{corollary}

\begin{proof}
We compare these walks with plane polygon classes by two injections.
Given a walk, add the four-edge cap
\[
 P_{0,0},\ R_{0,-1},\ P_{1,-1},\ R_{1,-1},\ P_{1,0}.
\]
The cap has no interior vertex in the half-plane, so it produces a simple polygon
whose unique lowest vertex is $P_{1,-1}$. That vertex fixes the
translation representative and makes the map injective into unrooted
classes.

For the reverse injection, take a plane polygon class and translate its
leftmost lowest vertex to $P_{0,0}$. A lowest vertex cannot be an $R$,
since such a vertex would use at least one descending bond. The chosen
$P$ vertex therefore uses both ascending diagonal bonds. Toggle the
hexagonal face whose upper path is
$P_{0,0},R_{0,0},P_{1,0}$ and whose lower path is the cap above:
that is, take symmetric edge difference with its boundary.

The polygon's overlap with this face consists of the first upper bond
and possibly the second. If $P_{1,0}$ is occupied, it too is lowest and
must use its bond to $R_{0,0}$. The three internal cap vertices lie below
the polygon. Thus the overlap is one connected
path containing all occupied face vertices. Toggling replaces that path
by the rest of the face and yields a simple capped polygon. Removing
the cap gives the required half-plane walk. The inverse adds the cap
and toggles the same fixed face, recovering the normalized polygon.

Both injections change length and diameter by bounded amounts and
weights by bounded positive factors. They transfer both tails in
\eqref{c4:Tail}, with walk length in the weighted sum. Their ratio
proves the conditional mean, and subtracting tails gives the
power-separated window.
\end{proof}

\subsection{Stacking polygons inside a strip}

\begin{proposition}[Confinement and the convergence radius]
\label{prop:strip-radius}
Let $x_{\rm strip}(W)$ be the length-series convergence radius of SAWs
from $P_{0,0}$ in the infinite vertical strip $|x|\le W$. Then
\[
 0\le x_{\rm strip}(W)-x_*\le W^{-4/3+o(1)},
\]
where the upper bound allows arbitrary fixed positive exponent slack.
\end{proposition}

\begin{proof}
The lower inequality follows from plane criticality. Fix
$0<\delta<1$ and $0<\epsilon<\delta/2$, and use plane polygon
classes with diameter in $[R^{1-\epsilon},R]$. By \eqref{c4:Tail},
their critical mass $A_{\rm ann}$ is polynomially large, and their
mean length is $R^{(4/3)(1-\epsilon)+o(1)}$. Set
\[
 z=x_*\exp(R^{-4/3+\delta}).
\]
Jensen's inequality gives
\[
 \log\sum_{\gamma\text{ in the annulus}}z^{|\gamma|}
 \ge \log A_{\rm ann}+R^{\delta-4\epsilon/3+o(1)}.
\]
The power in the last term is positive. Normalize the leftmost lowest
tip to $P_{0,0}$ and retain a subfamily with one common index $R_{i,j}$
for its leftmost highest tip. There are only $O(R^2)$ possible indices,
so this subfamily still has $z$-mass at least $\exp(R^c)$ for some
$c>0$. Here $j\ge1$, since one cell row has no cycle.

Stack any number of independently chosen copies upward. Put the next
lowest tip at the $P$ vertex vertically above the preceding designated
highest tip, and reflect every second copy left-to-right about the
vertical line through its own lowest tip. Reflection preserves the
lattice. Use the image of the designated $R_{i,j}$ as the top tip even
when it is no longer leftmost. The horizontal tip displacements now
alternate in sign, so every copy lies within horizontal distance $2R$
of the original start column.

At each junction toggle the facing hexagonal face extending to the
right of the common tip column. The lower polygon uses both descending
diagonal bonds at its highest $R$ tip, and the upper polygon uses both
ascending diagonal bonds at its lowest $P$ tip. Each overlap with the
joining face is therefore a path of one or two edges containing every
occupied face vertex, exactly as in the preceding cap argument.
Toggling joins the two polygons into one simple polygon. The two
overlap paths have two to four edges in total, so each join changes the
sum of the lengths by at most two.

Since $j\ge1$, successive joining faces are disjoint. For each fixed
stack size their locations are known from $(i,j)$ and the reflection
parity. Toggling those same faces recovers all the vertically separated
input polygons, so the construction is injective. The output remains
within horizontal distance $2R+3$ of the start column. Delete a fixed
bond at $P_{0,0}$, untouched at the very bottom, to obtain a SAW
starting there.

For each stack size, the total output $z$-mass is bounded below by the
product of the subfamily masses times a fixed positive factor per join
and a final bounded edge-removal factor. It grows without bound with
the stack size. These are subsets of the strip SAWs, so the strip
length series diverges at $z$; disjointness of outputs from different
stack sizes is not needed. Thus $x_{\rm strip}(2R+3)\le z$.
Rescaling $R$ and taking $\delta$ arbitrarily small proves the claim.
\end{proof}
